\chapter{MARCO TEÓRICO}

En este capítulo se desarrollan los fundamentos conceptuales y técnicos que sustentan la presente investigación. Se presentan, en primer lugar, las \textbf{arquitecturas de los modelos de clasificación} evaluados: Regresión Logística, XGBoost, Red Neuronal Multicapa (MLP) y LSTM con Mecanismo de Atención Aditiva. A continuación, se definen las \textbf{métricas de clasificación} empleadas para medir el rendimiento de dichos modelos en condiciones normales de operación. Posteriormente, se introducen los fundamentos del \textbf{Aprendizaje Automático Adversarial}, incluyendo la taxonomía de ataques de evasión en escenarios de caja blanca y caja negra, los algoritmos de ataque utilizados en la experimentación: CaFA, HopSkipJump, BoundaryAttack y SquareAttack, así como las \textbf{métricas de robustez adversarial} que permiten cuantificar el impacto operacional de dichos ataques sobre los sistemas de detección de fraude financiero y lavado de activos.

% ============================================================
\section{Modelos de Machine Learning}
% ============================================================

La detección de crimen financiero, que abarca tanto el fraude con medios de pago como el lavado de activos, constituye un problema de clasificación binaria, donde el objetivo es diferenciar automáticamente entre transacciones legítimas (clase 0) y transacciones ilícitas (clase 1) a partir de sus características transaccionales.

\subsection{Modelos Clásicos}

\subsubsection{Regresión Logística}

La \textbf{Regresión Logística} es un modelo de clasificación lineal que estima la probabilidad de que una transacción pertenezca a la clase positiva (actividad ilícita) mediante la función sigmoide aplicada a una combinación lineal de las características:
$$
P(y=1 \mid x) = \frac{1}{1 + e^{-(\mathbf{w}^\top x + b)}}
$$
donde $\mathbf{w}$ son los pesos del modelo, $x$ el vector de características y $b$ el sesgo. La clasificación final se obtiene comparando esta probabilidad contra un umbral de decisión $\tau$:
$$
\hat{y} = \begin{cases} 1 & \text{si } P(y=1 \mid x) \geq \tau \\ 0 & \text{en caso contrario} \end{cases}
$$
Su ventaja principal es la interpretabilidad, pues los pesos $\mathbf{w}$ indican directamente la contribución de cada característica a la predicción. Se utiliza como modelo lineal de referencia para comparar el rendimiento y la robustez adversarial frente a modelos más complejos \cite{hosmer2013logistic}.

\subsubsection{XGBoost}

El \textbf{XGBoost} (\textit{Extreme Gradient Boosting}) es un algoritmo de aprendizaje supervisado basado en \textit{boosting}, donde los árboles de decisión se construyen de forma secuencial, siendo cada uno entrenado para corregir los errores del anterior. La predicción final es la suma acumulada de todos los árboles:
$$
\hat{y}_i = \sum_{k=1}^{K} f_k(x_i)
$$
donde $K$ es el número de árboles y $f_k$ es la función del $k$-ésimo árbol.

Su principal ventaja en la detección de crimen financiero es el parámetro \texttt{scale\_pos\_weight}, que compensa el desbalance de clases asignando mayor peso a las transacciones ilícitas durante el entrenamiento. Esto lo convierte en el modelo de mayor rendimiento en datos tabulares desbalanceados, como los que caracterizan los escenarios reales de fraude y lavado de activos \cite{chen2016xgboost}.

\subsubsection{Red Neuronal Multicapa (MLP)}

La \textbf{Red Neuronal Multicapa} (\textit{Multi-Layer Perceptron}, MLP) es una arquitectura de aprendizaje profundo compuesta por una capa de entrada, una o más capas ocultas y una capa de salida. Cada neurona aplica una transformación no lineal mediante una función de activación $\sigma$, lo que permite al modelo capturar relaciones complejas entre las características de entrada:
$$
h^{(l)} = \sigma\left(W^{(l)} h^{(l-1)} + b^{(l)}\right)
$$
donde $h^{(l)}$ es la salida de la capa $l$, $W^{(l)}$ la matriz de pesos y $b^{(l)}$ el vector de sesgos. La capa final aplica una función \textit{softmax} para producir probabilidades de clase.

A diferencia de la Regresión Logística, el MLP puede aprender fronteras de decisión no lineales, lo que en principio lo hace más expresivo para datos tabulares complejos. Sin embargo, en escenarios con desbalance extremo de clases, su entrenamiento mediante descenso de gradiente estocástico puede colapsar hacia la predicción de la clase mayoritaria, resultando en un Recall nulo sobre la clase minoritaria \cite{rumelhart1986learning}.


% ============================================================
\section{Modelos de Redes Secuenciales y Atención}
% ============================================================

Las arquitecturas de \textit{Deep Learning} son relevantes en la detección de crimen financiero debido a su capacidad para modelar relaciones no lineales complejas entre variables transaccionales. En particular, los modelos recurrentes permiten capturar dependencias temporales en historiales de transacciones, lo que resulta útil para identificar patrones de comportamiento ilícito que se manifiestan a lo largo del tiempo, como transferencias estructuradas para evadir umbrales de reporte o cadenas de transacciones asociadas al lavado de activos.

\subsection{Long Short-Term Memory (LSTM)}

La \textit{Long Short-Term Memory} (LSTM) es un tipo de red neuronal recurrente diseñada para trabajar con datos secuenciales. A diferencia de las redes neuronales tradicionales, la LSTM puede recordar información por periodos prolongados gracias a un mecanismo interno de memoria. Su arquitectura incluye una celda de memoria y tres compuertas principales: la compuerta de olvido decide qué información anterior se descarta, la de entrada agrega nueva información relevante, y la de salida regula qué parte de la memoria influye en la siguiente predicción \cite{Benchaji2021}.

\begin{figure}[H]
\centering
\includegraphics[width=1\linewidth]{images/LSTM.png}
\caption{Estructura de la unidad LSTM \cite{Benchaji2021}}
\label{fig:LSTM_Standar}
\end{figure}

\subsection{Arquitectura LSTM con Mecanismo de Atención}

Si bien la LSTM es efectiva para modelar dependencias a largo plazo, puede presentar dificultades al manejar secuencias muy extensas. Para resolver esta limitación, se introduce el \textbf{Mecanismo de Atención}, que permite a la red ponderar selectivamente el historial de transacciones, enfocándose en aquellas más relevantes para la predicción actual. El mecanismo se aplica sobre la secuencia de estados ocultos $\mathbf{h}$ generados por la capa LSTM y se define en tres etapas \cite{Benchaji2021}:

\begin{enumerate}
    \item \textbf{Puntuación de alineación}: mide la relevancia del estado oculto actual $h_j$ respecto al estado previo $S_{t-1}$:
    \begin{equation}
        e_{tj} = a(S_{t-1}, h_{j})
    \end{equation}

    \item \textbf{Pesos de atención}: las puntuaciones se normalizan con \textit{Softmax}:
    \begin{equation}
        \alpha_{tj} = \frac{\exp(e_{tj})}{\sum_{k=1}^{n} \exp(e_{tk})}
    \end{equation}

    \item \textbf{Vector de contexto}: suma ponderada de los estados ocultos:
    \begin{equation}
        \mathbf{c}_t = \sum_{j=1}^{n} \alpha_{tj} \mathbf{h}_{j}
    \end{equation}
\end{enumerate}

El vector de contexto $\mathbf{c}_t$ se concatena con el último estado oculto y se introduce en la capa de clasificación final \cite{Benchaji2021}.

\begin{figure}[H]
\centering
\includegraphics[width=1\linewidth]{images/LSTM-Attention.png}
\caption{LSTM con Mecanismo de Atención \cite{Benchaji2021}}
\label{fig:lstm-attention}
\end{figure}


% ============================================================
\section{Métricas para los Modelos de Clasificación}
% ============================================================

\begin{enumerate}
    \item \textbf{\textit{Recall}}:
    Es la proporción de verdaderos positivos correctamente identificados por el modelo respecto al total de casos positivos reales. Evalúa la capacidad del modelo para detectar todas las transacciones ilícitas existentes en el conjunto de datos.
    \[
    \text{Recall} = \frac{TP}{TP + FN}
    \]
    Un valor alto de \textit{Recall} indica que el modelo logra capturar la mayoría de las transacciones ilícitas, reduciendo la cantidad de falsos negativos.

    \item \textbf{\textit{Precision}}:
    Mide la proporción de verdaderos positivos entre todas las predicciones positivas realizadas por el modelo. Indica qué fracción de las alertas generadas corresponden a transacciones ilícitas reales.
    \[
    \text{Precision} = \frac{TP}{TP + FP}
    \]
    Un valor alto de \textit{Precision} significa que la mayoría de las alertas generadas son correctas, reduciendo los falsos positivos.

    \item \textbf{F1-Score}:
    Es la media armónica entre la \textit{Precision} y el \textit{Recall}. Proporciona una medida equilibrada del rendimiento, especialmente útil cuando existe desbalance entre clases, ya que penaliza modelos que sacrifican una métrica en favor de la otra.
    \[
    \text{F1-Score} = 2 \cdot \frac{\text{Precision} \cdot \text{Recall}}{\text{Precision} + \text{Recall}}
    \]

    \item \textbf{AUC-ROC} \textit{(Area Under the Receiver Operating Characteristic Curve)}:
    Mide la capacidad del modelo para distinguir entre transacciones ilícitas y legítimas independientemente del umbral de decisión. Se define como el área bajo la curva que relaciona la Tasa de Verdaderos Positivos (TPR) contra la Tasa de Falsos Positivos (FPR) a distintos umbrales:
    \[
    \text{AUC-ROC} = \int_0^1 TPR \, d(FPR)
    \]
    Un valor de 1.0 indica separación perfecta entre clases; un valor de 0.5 equivale a una clasificación aleatoria. A diferencia del F1-Score, el AUC-ROC no se ve afectada por el desbalance de clases, lo que la convierte en la métrica principal de comparación entre modelos en este trabajo.
\end{enumerate}


% ============================================================
\section{Aprendizaje Automático Adversarial (AML)}
% ============================================================

El Aprendizaje Automático Adversarial (Adversarial Machine Learning, AML) examina las vulnerabilidades de los modelos de aprendizaje automático frente a perturbaciones intencionadas en los datos de entrada \cite{biggio2018wild}. En este contexto, la \textit{robustez} refleja la capacidad de un modelo para mantener un rendimiento estable ante manipulaciones adversariales \cite{fawzi_robustness_2016}: un modelo robusto no solo debe generalizar correctamente ante datos no vistos, sino también resistir modificaciones deliberadas en sus entradas.

Esta interacción se formaliza como un problema de optimización. El ataque adversarial busca la perturbación $\delta$ que maximiza la función de pérdida del modelo:
\[
A_i(\theta, \mathcal{B}) = \max_{\|\delta_i\| \leq \epsilon} L\big(f_\theta(x_i + \delta_i), y_i\big),
\]
donde $L$ es la función de pérdida, $f_\theta$ el modelo parametrizado por $\theta$, y $\mathcal{B} = \{\delta \in \mathbb{R}^d : \|\delta\| \leq \epsilon\}$ el conjunto de perturbaciones permitidas.

El impacto global se evalúa mediante el \textit{riesgo adversarial empírico}, promedio de las pérdidas sobre los $n$ ejemplos:
\[
R_{P_n}(\theta, \mathcal{B}) = \frac{1}{n} \sum_{i=1}^{n} A_i(\theta, \mathcal{B}).
\]

Como medida de defensa, el \textit{entrenamiento robusto} optimiza los parámetros del modelo minimizando este riesgo \cite{goodfellow2015explaining}:
\[
\theta^{\star} = \arg\min_{\theta} R_{P_n}(\theta, \mathcal{B}).
\]

En la detección de crimen financiero, que comprende tanto el fraude con medios de pago como el lavado de activos, el AML es especialmente relevante, ya que un atacante puede modificar características como el monto, la ubicación o la temporalidad de una transacción para imitar el comportamiento de una operación legítima y evadir la detección. El presente trabajo se enmarca en este contexto, evaluando la vulnerabilidad de distintos clasificadores financieros frente a ataques de evasión y analizando el costo mínimo de perturbación necesario para comprometer su capacidad de detección.


\subsection{Ataques Adversariales de Evasión}

Los ataques adversariales pueden clasificarse según la etapa del ciclo de vida del modelo en la que se producen. Entre las categorías más comunes se encuentran los \textit{ataques de envenenamiento}, que corrompen el conjunto de entrenamiento para modificar el comportamiento futuro del modelo; los \textit{ataques de extracción}, que buscan replicar o robar la funcionalidad del modelo; y los \textit{ataques de inferencia}, cuyo propósito es revelar información sensible de los datos de entrenamiento \cite{huang2011adversarial}. Sin embargo, la presente investigación se enfoca en los \textit{ataques de evasión}, por ser los más frecuentes y realistas en escenarios de crimen financiero: el atacante modifica una transacción ilícita en el momento de la inferencia para que el modelo la clasifique como legítima, sin alterar el modelo ni los datos de entrenamiento.

Los ataques de evasión se ejecutan durante la fase de \textit{inferencia}, cuando el modelo ya ha sido entrenado y desplegado. El adversario busca generar una entrada modificada que logre evadir la detección sin alterar los parámetros internos del modelo. Dicha entrada se conoce como \textit{ejemplo adversarial} y surge de aplicar una perturbación $\delta$ a una transacción ilícita $x$, generando una nueva instancia $x' = x + \delta$ que produce una clasificación errónea:
\[
f_\theta(x') \neq f_\theta(x), \quad \text{con} \quad \|\delta\|_p \leq \epsilon,
\]
donde $\epsilon$ controla la magnitud de la perturbación y $\|\cdot\|_p$ es la norma utilizada para medirla. Las modificaciones suelen ser de baja amplitud, coherentes desde el punto de vista financiero, pero suficientes para alterar la decisión del modelo.

La efectividad de este tipo de ataque depende del nivel de acceso que el adversario tenga sobre el modelo, lo que da lugar a dos escenarios principales: \textit{caja blanca} (\textit{white-box}), con acceso total al modelo, y \textit{caja negra} (\textit{black-box}), donde solo se observan las salidas \cite{huang2011adversarial}.

\subsubsection{Ataques de Caja Blanca}

En un escenario de caja blanca, el adversario tiene acceso completo al interior del modelo: conoce su arquitectura, sus parámetros y cómo calcula el error. Con esta información puede calcular gradientes exactos y modificar una transacción de forma dirigida para que el modelo la clasifique incorrectamente \cite{madry2018towards}.

En el presente trabajo se utiliza \textbf{CaFA} (\textit{Categorical Feature Attack}) \cite{cartella2021adversarial} para representar este escenario. CaFA fue diseñado específicamente para datos tabulares financieros, donde no basta con modificar números arbitrariamente, ya que las perturbaciones deben tener sentido en el mundo real. Por ello, tras cada modificación, CaFA verifica que el ejemplo resultante siga siendo válido: que los montos no sean negativos, que las variables categóricas tomen solo valores permitidos y que los campos sean coherentes entre sí. Esto lo diferencia de otros ataques basados en gradientes que fueron diseñados originalmente para imágenes y no consideran las restricciones propias del dominio financiero.

\subsubsection{Ataques de Caja Negra}

El adversario no tiene acceso a la arquitectura ni a los parámetros del modelo, y solo puede interactuar con él mediante consultas y respuestas de salida \cite{shokri2017membership}. Este escenario es el más realista en entornos de producción, donde los sistemas de clasificación operan como servicios cerrados.

Según la información devuelta por el modelo, estos ataques se clasifican en:
\begin{itemize}
    \item \textit{Score-based}: el adversario obtiene probabilidades por clase y las utiliza para guiar la búsqueda de perturbaciones efectivas.
    \item \textit{Decision-based}: el sistema solo devuelve etiquetas, por lo que el atacante infiere la frontera de decisión mediante muestreo iterativo.
\end{itemize}

En el presente trabajo se emplean tres ataques de caja negra: \textbf{HopSkipJump} y \textbf{BoundaryAttack}, ambos de tipo \textit{decision-based}, y \textbf{SquareAttack}, de tipo \textit{score-based}. Este escenario es especialmente relevante en la detección de crimen financiero, ya que los sistemas de detección operan típicamente como servicios cerrados donde el atacante solo puede observar la respuesta del modelo y ajustar sus transacciones en consecuencia.


\subsection{Algoritmos de Ataque de Caja Negra}

El \textbf{HopSkipJump Attack} (HSJ) \cite{chen2019hopskipjump} es un ataque \textit{decision-based} que solo requiere las etiquetas de salida del modelo. Su objetivo es minimizar la distancia entre el ejemplo original y el ejemplo adversarial bajo la restricción de cambio de etiqueta:
\[
\delta^\star = \arg\min_{\delta} \|\delta\|_p \quad \text{s.a.} \quad f(\Pi_{\mathcal{C}}(x+\delta)) \neq f(x).
\]
HSJ estima direcciones hacia la frontera de decisión mediante muestreo estocástico y búsquedas binarias. En datos tabulares, la proyección $\Pi_{\mathcal{C}}$ garantiza que las características modificadas permanezcan dentro de valores financieramente plausibles.

El \textbf{Boundary Attack} \cite{brendel2018decision} también es \textit{decision-based} y explora la frontera de decisión partiendo de un punto adversarial inicial $x_0$ tal que $f(x_0) \neq f(x)$. En cada iteración genera propuestas de la forma:
\[
\tilde{x} = x_t + \gamma \frac{z}{\|z\|}, \qquad z \sim \mathcal{N}(0,I),
\]
que se proyectan hacia el ejemplo original obteniendo $x' = \Pi_{\mathcal{C}}\big(\Pi_{\mathrm{line}}(x,\tilde{x})\big)$. Si $f(x') \neq f(x)$, el punto se acepta y se ajusta $\gamma$ para reducir progresivamente la distancia al original.

El \textbf{Square Attack} \cite{andriushchenko2020squareattackqueryefficientblackbox} es un ataque \textit{score-based} que requiere las probabilidades de salida del modelo. Perturba iterativamente subconjuntos aleatorios de características, aceptando únicamente las modificaciones que incrementan la pérdida del clasificador:
\[
x' = \Pi_{\mathcal{C}}\big(x + \delta_S\big),
\]
donde $\delta_S$ actúa solo sobre un subconjunto $S$ de características con magnitud acotada por $\|\delta_S\|_\infty \leq \varepsilon$. Este enfoque resulta más eficiente en consultas que los métodos basados en estimación de gradiente y ha demostrado alta efectividad en datos tabulares con perturbaciones de bajo costo $L_0$.


\subsection{Métricas de Evaluación de Robustez Adversarial}

Para evaluar la efectividad de los ataques y la robustez de los modelos, se emplean las siguientes métricas experimentales:

\begin{itemize}
    \item \textbf{Tasa de Evasión}: Proporción de ejemplos adversariales que logran ser clasificados como legítimos por el modelo objetivo tras la perturbación. Mide directamente si el ataque tuvo éxito:
    \[
    \text{Tasa de Evasión} = \frac{\text{muestras mal clasificadas post-ataque}}{\text{total de muestras atacadas}}.
    \]

    \item \textbf{$\Delta$FNR (Incremento en la Tasa de Falsos Negativos)}: Mide el daño operacional real causado por el ataque, cuantificando cuántas transacciones ilícitas adicionales deja de detectar el modelo:
    \[
    \Delta\text{FNR} = \text{FNR}_{\text{post-ataque}} - \text{FNR}_{\text{pre-ataque}},
    \]
    donde $\text{FNR} = \dfrac{FN}{FN + TP}$, siendo $FN$ los falsos negativos y $TP$ los verdaderos positivos. Un $\Delta\text{FNR} = 1.0$ indica que todas las transacciones ilícitas detectadas antes del ataque ahora escapan al detector.

    \item \textbf{Recall post-ataque}: Proporción de transacciones ilícitas que el modelo sigue detectando correctamente como verdaderos positivos después del ataque:
    \[
    \text{Recall}_{\text{post}} = \frac{TP_{\text{post}}}{TP_{\text{post}} + FN_{\text{post}}}.
    \]

    \item \textbf{Norma $L_0$ (Costo de Perturbación)}: Número de características que el atacante modificó para generar el ejemplo adversarial. Un $L_0$ bajo indica un ataque sigiloso, ya que pocas modificaciones bastan para engañar al modelo:
    \[
    L_0 = \left|\{ i : x_i \neq x'_i \}\right|,
    \]
    donde $x$ es la muestra original y $x'$ es la muestra adversarial. Se reportan la media, mínimo, máximo y desviación estándar sobre todas las evasiones exitosas.

    \item \textbf{$L_\infty$ estandarizada (Magnitud de la Perturbación)}:
    Magnitud máxima de cambio aplicada a cualquier característica, normalizada por el rango válido de esa variable. Un valor cercano a 0 indica que las modificaciones son imperceptibles dentro del rango normal de operación, dificultando su detección por reglas de validación.
\end{itemize}
